\section{Review of Related Literature}\label{review-of-related-literature}

\subsection{Review Basis and Organization}\label{review-basis-and-organization}
This chapter synthesizes two complementary evidence streams. The first
concerns urban drainage monitoring, blockage, graph-based vulnerability
assessment, critical-component analysis, and hydraulic or network-level
prioritization. This literature establishes the practical origin of the
research problem. The second concerns foundational and modern
propagation models across influence diffusion, epidemics, cascading
failure, cybersecurity, finance, transportation, and infrastructure.
This literature establishes the mechanisms and state semantics that can
inform a general graph-risk architecture.

The review follows a two-stage synthesis. It first establishes the
domain-specific problem through drainage monitoring and graph-based
prioritization. It then situates that problem within the broader
literature on propagation-state semantics, topology, transmission,
susceptibility, recursion, aggregation, dependence, model complexity,
and ranking. The final sections compare reusable architectures and
domain-specific parameterization before synthesizing the research
problems and the architectural gap addressed by the study.

\subsection{Urban Drainage Networks as Interconnected
Systems}\label{urban-drainage-networks-as-interconnected-systems}

Urban drainage networks consist of connected conduits, junctions,
inlets, screens, and other components whose behavior is influenced by
flow direction, network structure, hydraulic condition, and local
impairment. Blockage, clogging, siltation, structural defects, or
reduced effective capacity can influence system behavior beyond the
initially affected location. Numerical and network-based drainage
studies show that the location and severity of impairment can alter
hydraulic response, vulnerability, and the importance of connected
components \citep{simone2022,simone2023,reyes2020,dastgir2023,dastgir2024,meijer2023,zhang2025}.

This interconnected structure creates a practical prioritization
problem. Inspection and maintenance resources are finite, so not every
component can receive equal attention at the same time. Existing
drainage studies explicitly frame critical-component identification as
useful for targeted management, maintenance, monitoring, resilience, or
dredging decisions \citep{simone2022}, \citep{simone2023}, \citep{dastgir2023,dastgir2024,meijer2023,zhang2025}. The practical
decision problem is therefore not only to identify which components are
locally impaired but also to determine which components deserve greater
attention when the surrounding network is considered.

This drainage-side evidence exposes a distinction between observation
and prioritization. A local
blockage or defect can be observed directly, but a maintenance decision
often requires a broader question: whether that component remains
important after the connected network, downstream pathways, relation
conditions, and converging influences are considered. This distinction
between local condition and network-context priority recurs in studies
that combine inspection, topology, capacity, and hydraulic relevance
\citep{simone2022,simone2023,dastgir2023,meijer2023,zhang2025}.

\subsection{Local Monitoring and Component-Level Condition
Assessment}\label{local-monitoring-and-component-level-condition-assessment}

Computer-vision, sensor, and inspection studies have substantially
improved the ability to identify local drainage condition. \citet{iqbal2021}, for example, formulate culvert blockage as an
image-classification problem and evaluate several convolutional neural
networks for classifying blockage from imagery. \citet{vandaele2024}
similarly develop camera-based methods for automatically detecting
blockage at trash screens for maintenance and flood-management use.
These approaches demonstrate that local blockage information can be
generated automatically and repeatedly at monitored assets.

The outputs of these methods are valuable because they provide repeatable
local condition or disturbance signals. However, a local score by itself
does not determine network-context priority: it describes the monitored
asset rather than the recursively accumulated state of the wider connected
system \citep{iqbal2021,vandaele2024,simone2022}.

This distinction matters when two components have similar observed local
conditions but occupy different positions in the network. One may lie on
a relation with high operating susceptibility or feed several important
downstream paths, while another may have little propagated consequence
under the same network configuration. Local assessment is therefore
necessary for condition awareness, but it does not by itself answer which
nodes warrant greater priority after network position and operating
conditions are considered \citep{simone2022,dastgir2023,zhang2025}.

The limitation should not be interpreted as a weakness of computer
vision or sensing. Those methods solve the detection problem they were
designed to solve. The unresolved issue is the interface between local
condition information and network-level prioritization. This motivates
the need for a propagation layer that can accept a local disturbance as
input while separately representing network structure and
relation-specific response \citep{simone2022,dastgir2023,meijer2023}.

\subsection{Graph-Based Drainage Analysis and Structural
Prioritization}\label{graph-based-drainage-analysis-and-structural-prioritization}

Drainage research has increasingly used graph theory and complex-network
analysis to identify critical or vulnerable components. \citet{simone2022} examine standard and tailored centrality metrics for
vulnerability assessment, monitoring design, and contaminant or
pathogenic spread. Their analysis shows that standard degree, harmonic,
betweenness, and edge-betweenness measures do not always reproduce
hydraulically meaningful importance, which motivates tailored metrics
that incorporate intrinsic relevance and flow direction.

\citet{simone2023} extends this idea using relevance-based centrality for
critical-pipe identification. The method ranks nodes and pipes by
combining network connectivity with intrinsic relevance rather than
using plain topological centrality alone. This is significant for the
present study because it demonstrates that even within a graph-only
framework, useful prioritization can require information beyond raw
connectivity.

\citet{dastgir2023} propose a graph method for critical-pipe
analysis based on runoff edge betweenness and capacity-related
information. Their method is compared with hydrodynamic simulation and
is intended to retain useful criticality information with lower
computational and data requirements. The study also reports that
graph-method accuracy decreases as network complexity, loop degree, and
hydraulic stress increase, illustrating the tradeoff between lightweight
structural approximation and richer physical behavior.

\citet{dastgir2024} subsequently examine multiple-pipe failures in
stormwater networks using graph-theoretic representations. Their work is
particularly relevant to the present study because it shows that the
consequences and prioritization of interacting failures cannot always be
inferred from isolated single-component analysis. It therefore supports
the explicit evaluation of multi-source configurations and converging
effects while also reinforcing the need to distinguish graph-based
screening from full hydraulic consequence prediction.

\citet{meijer2023} develop a graph-based weakest-link method for
urban drainage systems using flow-path analysis instead of repeated
hydrodynamic simulation. Their work focuses on storage and
discharge-capacity relationships and demonstrates that graph-based
screening can make long rainfall-series analysis practically feasible.
This provides further evidence that drainage research already seeks
computationally efficient prioritization and vulnerability tools rather
than relying exclusively on full hydraulic simulation.

\citet{zhang2025} analyze siltation risk using complex-network
concepts, centrality, pipeline risk factors, and edge vulnerability.
Their study identifies critical nodes and high-risk propagation
pathways, explicitly linking siltation-chain structure with maintenance
and dredging prioritization. This is especially relevant because it
shows that drainage risk prioritization may involve both structural
importance and domain-specific condition information rather than a
single metric.

Taken together, these studies establish two important points. First,
prioritization is already a legitimate drainage-network task; the
novelty of the present study is therefore not the act of ranking nodes
or pipes. Second, the literature repeatedly augments topology with
runoff, hydraulic relevance, storage/discharge capacity, edge
vulnerability, siltation condition, or other operational information.
Structural position and operating response are therefore related but not
equivalent \citep{simone2022,simone2023,dastgir2023,meijer2023,zhang2025}.

\subsection{Propagation Models and State
Semantics}\label{propagation-models-and-state-semantics}

Broader network science provides several foundational propagation
paradigms. Independent Cascade and Linear Threshold models represent
influence through activation semantics \citep{kempe2003}. Epidemic models such as
SIR and SIS propagate compartment or infection states \citep{kermack1927,pastor2001,humphries2021}.
Cascading-failure models propagate physical load, capacity exceedance,
or component failure \citep{buldyrev2010,he2019,shuang2017}. Attack graphs represent feasible
attack paths and reachable compromise states, while Bayesian
attack-graph approaches introduce probabilistic compromise semantics
\citep{sheyner2002}, \citep{munoz2017}. Cyber-risk, finance, transportation, and
infrastructure models introduce additional states such as accumulated
risk, exposure, payment loss, systemic stress, or dynamic load
\citep{kavallieratos2021,faraji2026,zhou2023,wu2026,petrone2018}.

The main lesson is that the meaning of the propagated state determines
the valid update and aggregation rules. Binary activation answers
whether a node changes state. Probability estimates event likelihood
under a probabilistic model. Physical-load models track conserved or
redistributed quantities. These states are appropriate for their
intended questions, but they do not automatically answer the comparative
question of which node carries greater comparative network-context risk
under a common parameterization \citep{kempe2003,kermack1927,buldyrev2010,
sheyner2002,kavallieratos2021}.

The comparison indicates that any proposed propagation model must declare
whether its state represents activation, probability, physical quantity,
or a non-probabilistic index. Without that declaration, similarly bounded
values may be interpreted incorrectly and rankings from different model
families may appear comparable when they are not \citep{kempe2003,
kermack1927,munoz2017,petrone2018}.

Table~\ref{tab:model-families} compares the model families at the level
emphasized during consultation: their structure, propagated state,
update basis, analytical purpose, and limitation for a lightweight
comparative-risk task. The comparison shows why no single family can be
adopted merely because it uses a graph; the semantics of its state and
the assumptions of its update rule must match the research question.

\begingroup
\small
\begin{longtable}{@{}p{0.14\textwidth}p{0.15\textwidth}p{0.18\textwidth}p{0.17\textwidth}p{0.19\textwidth}@{}}
\caption{Comparison of propagation-model families relevant to comparative network-risk prioritization.}\label{tab:model-families}\\
\toprule
Model family & Typical state & Propagation or update basis & Principal analytical use & Limitation for the present problem \\
\midrule
\endfirsthead
\caption[]{Comparison of propagation-model families relevant to comparative network-risk prioritization (continued).}\\
\toprule
Model family & Typical state & Propagation or update basis & Principal analytical use & Limitation for the present problem \\
\midrule
\endhead
\bottomrule
\endfoot
Influence diffusion \citep{kempe2003} & Binary activation or threshold state & Edge influence probabilities or accumulated neighbor influence & Spread maximization and activation reach & Activation semantics do not directly represent a continuous comparative risk index. \\
Epidemic propagation \citep{kermack1927,pastor2001,humphries2021} & Susceptible, infected, and recovered compartments or infection probabilities & Contact structure, infection rate, and recovery rate & Disease spread and intervention analysis & Epidemiological compartments and rates require domain-specific assumptions not implied by network priority. \\
Cascading failure \citep{buldyrev2010,he2019,shuang2017} & Load, capacity, operational state, or failure state & Load redistribution, tolerance, and interdependence & Robustness and cascade consequence & Physical load conservation or failure thresholds are not generic comparative-risk semantics. \\
Attack graph and Bayesian attack graph \citep{sheyner2002,munoz2017} & Reachable compromise state or compromise probability & Feasible attack steps, conditional probabilities, or logical relations & Security-path and compromise analysis & Probabilistic outputs require calibrated conditional assumptions and dependence treatment. \\
Financial and infrastructure contagion \citep{faraji2026,zhou2023,wu2026} & Exposure, loss, stress, or service degradation & Dependency, exposure, vulnerability, and feedback & Systemic-risk and resilience assessment & Feedback and common exposure can require cyclic or dependence-aware dynamics. \\
Graph-based drainage screening \citep{simone2022,simone2023,dastgir2023,meijer2023,zhang2025} & Centrality, criticality, capacity, or vulnerability score & Topology combined with relevance, runoff, capacity, or siltation information & Monitoring, maintenance, and critical-component screening & Scores differ in meaning and may omit recursive current-state propagation. \\
Hydraulic simulation and learned surrogate \citep{reyes2020,dastgir2024,zhang2024} & Flow, depth, surcharge, flood volume, or learned hydraulic state & Governing physics or simulator-trained approximation & Physical consequence analysis and prediction & Greater data, calibration, or training requirements; outputs answer a different question from lightweight comparative risk. \\
\end{longtable}
\endgroup


\subsection{Topology, Reachability, and
Transmission}\label{topology-reachability-and-transmission}

Across propagation models, topology primarily determines which
interactions are possible. An edge $i\to j$ indicates that influence may move
from node i to node j, but the existence of the edge does not by itself
define the strength of influence. Influence-diffusion models, attack
graphs, cyber-risk models, and weighted infrastructure models commonly
distinguish the existence of a pathway from the weight, probability,
exposure, or transfer mechanism associated with that pathway
\citep{kempe2003,sheyner2002,kavallieratos2021,faraji2026,petrone2018}.

This distinction is particularly important for interpreting centrality.
Centrality uses global graph structure and can identify structurally
important nodes or edges, but it does not directly simulate the
propagation of a current disturbance under current relation-specific
conditions. \citet{simone2022} provide a drainage-specific example:
standard centrality metrics can rank structurally central elements
differently from hydraulically relevant elements, which motivated the
incorporation of relevance and direction.

Accordingly, the literature supports treating topology as reachability
while retaining transmission or influence as a conceptually separate
relation-level mechanism. Where heterogeneous influence cannot be
empirically justified, a neutral coefficient is more defensible than an
unsupported set of edge weights \citep{kempe2003,simone2022,dastgir2023}.

\subsection{Susceptibility, Capacity, Vulnerability, and Receiver
Response}\label{susceptibility-capacity-vulnerability-and-receiver-response}

A second distinction concerns transmission versus susceptibility.
Transmission describes how strongly an influence transfers along a
relation, whereas susceptibility describes how strongly a relation or
receiving context responds under the current condition. Similar
distinctions appear in models using capacity, vulnerability, tolerance,
resilience, exposure, or fragility \citep{buldyrev2010,he2019,
kavallieratos2021,faraji2026}.

Cascading-failure models demonstrate that the same incoming load can
produce different outcomes depending on component capacity or tolerance
\citep{buldyrev2010,he2019,shuang2017}. Cyber and infrastructure models likewise separate
dependency or exposure strength from vulnerability or resilience
\citep{kavallieratos2021,faraji2026,zhou2023}. Drainage studies that incorporate hydraulic
relevance, runoff, capacity, storage/discharge relationships, siltation
condition, or edge vulnerability similarly show that structural
connectivity alone is not sufficient to characterize operating response
\citep{simone2022,simone2023,reyes2020,dastgir2023,dastgir2024,meijer2023,zhang2025}.

These patterns support keeping relation-specific susceptibility distinct
from transmission rather than collapsing both concepts into one
unexplained coefficient. In a drainage application, susceptibility may be
linked to load relative to effective capacity, whereas heterogeneous
transmission should remain neutral unless empirical evidence justifies its
calibration \citep{dastgir2023,meijer2023,zhang2025}.

\subsection{Recursive Propagation and Current-State
Updating}\label{recursive-propagation-and-current-state-updating}

Many propagation systems are recursive: later effects depend on the
state produced by earlier propagation rather than repeatedly referencing
only the original source. Epidemic models update infection state over
successive interactions, cascading-failure models redistribute load
after each failure, financial-contagion models update exposure or loss
as conditions change, and hydraulic models update system state in
response to preceding conditions \citep{kermack1927,buldyrev2010,
he2019,faraji2026,petrone2018}.

This pattern supports current-state recursion: for a chain
$A\to B\to C$, the contribution from B to C should be based on the current
modeled state at B when the intended semantics permit inherited and local
influence to accumulate. Conceptually, the node propagates what it
currently carries rather than only the disturbance that originally
entered the network \citep{kermack1927,buldyrev2010,faraji2026}.

Current-state recursion is central to comparative network context. If
the model repeatedly used only the original local disturbance,
intermediate accumulation would be lost and downstream ranking would not
reflect the evolving state of the network \citep{kempe2003,buldyrev2010,
petrone2018}.

\subsection{Multi-Source Aggregation and Bounded
State}\label{multi-source-aggregation-and-bounded-state}

When several upstream relations converge on the same node, their
contributions require aggregation. Existing model families use different
operators according to state semantics: physical loads may be summed or
redistributed, threshold models accumulate influence until a threshold
is crossed, and probabilities may be combined through Bayesian or
noisy-OR structures when the required assumptions are satisfied
\citep{kempe2003,munoz2017,buldyrev2010}.

For a bounded non-probabilistic comparative risk state, simple addition
creates a practical difficulty because multiple contributions can exceed
the declared {[}0,1{]} range. Clipping can restore the range but introduces
an abrupt ceiling that loses information about how additional
contributions accumulate. Multiplicative-complement aggregation is a
candidate bounded, saturating operator, although its resemblance to
noisy-OR algebra requires an explicit warning: it has a probabilistic
interpretation only when the inputs and dependence assumptions justify
one \citep{munoz2017}. The selected operator must therefore follow the
declared state semantics rather than visual similarity to an established
probability formula.

\subsection{Dependence, Shared Paths, and
Reconvergence}\label{dependence-shared-paths-and-reconvergence}

Bounded aggregation does not solve every multi-source problem. When
contributions share a common origin and travel through multiple paths
that later reconverge, incoming contributions may not be independent
descriptions of distinct risk sources. Similar dependence problems
appear in financial networks through common exposures, in infrastructure
networks through shared dependencies, and in graph systems with
overlapping paths \citep{buldyrev2010,faraji2026,zhou2023}.

The literature therefore supports treating local aggregation and
dependence as separate mechanisms. Bounding incoming contributions does
not by itself correct shared-origin or common-cause dependence;
reconvergence should be evaluated as an explicit limitation rather than
hidden inside the aggregation rule \citep{buldyrev2010,faraji2026}.

This distinction is important for interpretability. A bounded output
does not automatically mean that shared-origin information has been
counted correctly. Dependence-aware correction remains a separate
extension problem \citep{munoz2017,faraji2026}.

\subsection{Lightweight Graph Models, Learned Surrogates, and
High-Fidelity
Simulation}\label{lightweight-graph-models-learned-surrogates-and-high-fidelity-simulation}

The drainage literature contains a spectrum of model complexity. Local
scores and graph indicators provide relatively inexpensive screening;
graph-hydrodynamic hybrids combine structural reasoning with physical
analysis; hydraulic simulators explicitly represent domain dynamics; and
graph-neural-network surrogates learn simulator-derived hydraulic
behavior for faster inference \citep{simone2022,dastgir2023,meijer2023,
zhang2024}.

\citet{dastgir2023} explicitly compare a graph-based critical-pipe
method with hydrodynamic modeling and report substantially lower data
and computational requirements for the graph approach, while also
documenting accuracy loss as complexity and stress increase. Simone et
al. \citep{simone2022} likewise position tailored complex-network metrics as
complementary preliminary tools when full hydraulic simulation is more
complex, data demanding, or computationally expensive.

At the learned-surrogate end of the spectrum, \citet{zhang2024}
develop a graph-neural-network surrogate for urban-drainage hydraulic
prediction. The model is trained against physics-based flow-routing
outputs and predicts hydraulic variables for real-time use.
Physics-guided mechanisms are included to improve prediction
consistency. Such a surrogate seeks to emulate hydraulic state transitions,
which is a different objective from transparent graph screening or
comparative index construction \citep{zhang2024}.

The literature therefore does not imply that lightweight models are
universally superior. High-fidelity models provide richer physical
detail, while learned surrogates can approximate that detail after
training. Lightweight graph models instead trade physical completeness
for lower data requirements, simpler computation, and direct
interpretability. The appropriate model depends on the decision
question \citep{dastgir2023,simone2022,zhang2024}.

\subsection{Existing Ranking and Prioritization
Approaches}\label{existing-ranking-and-prioritization-approaches}

Existing drainage and network studies rank components using different
criteria. Local monitoring may rank according to observed blockage or
defect severity. Centrality approaches rank according to structural
importance. Weakest-link and vulnerability methods incorporate
combinations of topology and domain-specific information. Failure
studies rank components or combinations according to modeled
consequence. Siltation studies identify high-risk nodes and propagation
pathways. Hydraulic approaches may prioritize locations according to
depth, surcharge, flood volume, resilience, or performance loss
\citep{iqbal2021,vandaele2024,simone2022,dastgir2023,dastgir2024,
meijer2023,zhang2025}.

The evidence shows that ranking is not a single modeling operation with
a universal meaning. A centrality rank, blockage-severity rank,
weakest-link rank, hydraulic-impact rank, and propagated-risk rank
answer different questions. The semantic basis of the score must
therefore be explicit before rankings can be compared
\citep{simone2022,dastgir2023,meijer2023,zhang2025}.


% \begin{quote}\textit{Table migration note: this comparison table was preserved separately in tables\_pending.tex and will be rebuilt cleanly in LaTeX before final submission.}\end{quote}

{\small
\renewcommand{\arraystretch}{1.15}
\begin{longtable}{@{}>{\raggedright\arraybackslash}p{0.20\dimexpr\linewidth-6\tabcolsep\relax}
                     >{\raggedright\arraybackslash}p{0.27\dimexpr\linewidth-6\tabcolsep\relax}
                     >{\raggedright\arraybackslash}p{0.27\dimexpr\linewidth-6\tabcolsep\relax}
                     >{\raggedright\arraybackslash}p{0.26\dimexpr\linewidth-6\tabcolsep\relax}@{}}
\caption{Comparison of ranking and prioritization approaches}\label{tab:ranking-approaches}\\
\toprule
\textbf{Approach} & \textbf{Primary ranking basis} & \textbf{What the ranking answers} & \textbf{Limitation relative to the present target} \\
\midrule
\endfirsthead
\toprule
\textbf{Approach} & \textbf{Primary ranking basis} & \textbf{What the ranking answers} & \textbf{Limitation relative to the present target} \\
\midrule
\endhead
\bottomrule
\endlastfoot

Local condition / detection~\citep{iqbal2021,vandaele2024,patil2023} &
Observed blockage, defect, debris, or local score &
Which asset appears more locally impaired? &
Does not inherently incorporate recursively propagated network context \\
\addlinespace

Centrality / structural criticality~\citep{simone2022,simone2023,reyes2020} &
Graph position, paths, relevance, or structural interaction &
Which node or edge is more structurally important? &
Does not by itself represent propagation of a current disturbance under current relation conditions \\
\addlinespace

Weakest-link / vulnerability methods~\citep{dastgir2023,meijer2023} &
Structural information combined with capacity, storage, discharge, redundancy, or other domain variables &
Which component most constrains or weakens system performance? &
Produces specialized vulnerability or criticality semantics rather than a generic bounded recursive risk state \\
\addlinespace

Siltation-chain risk~\citep{zhang2025} &
Pipeline risk, node importance, edge vulnerability, and propagation path &
Which nodes and pathways are high priority for siltation management? &
Domain-specific siltation semantics and parameterization \\
\addlinespace

Hydraulic consequence analysis~\citep{dastgir2024,zhang2024} &
Simulated depth, surcharge, flood volume, performance loss, or resilience &
Which failures or locations cause greater physical consequence? &
High-fidelity but strongly domain-specific and more data/computation intensive \\

\end{longtable}
}



The literature comparison also separates score generation from the
decision-support use of a score. A method may first compute local,
structural, hydraulic, or propagated values and then order components for
inspection or maintenance; the resulting ranking inherits the meaning
and limitations of the underlying score \citep{simone2022,dastgir2023,
meijer2023,zhang2025}.

\subsubsection{Drainage Prioritization Evidence
Matrix}\label{drainage-prioritization-evidence-matrix}

The following study-level comparison makes the progression from local
condition assessment to structural, hybrid, risk-path, and learned
hydraulic approaches explicit.


% \begin{quote}\textit{Table migration note: this comparison table was preserved separately in tables\_pending.tex and will be rebuilt cleanly in LaTeX before final submission.}\end{quote}

{\small
\renewcommand{\arraystretch}{1.15}
\begin{longtable}{@{}>{\raggedright\arraybackslash}p{0.14\dimexpr\linewidth-8\tabcolsep\relax}
                     >{\raggedright\arraybackslash}p{0.17\dimexpr\linewidth-8\tabcolsep\relax}
                     >{\raggedright\arraybackslash}p{0.23\dimexpr\linewidth-8\tabcolsep\relax}
                     >{\raggedright\arraybackslash}p{0.22\dimexpr\linewidth-8\tabcolsep\relax}
                     >{\raggedright\arraybackslash}p{0.24\dimexpr\linewidth-8\tabcolsep\relax}@{}}
\caption{Drainage prioritization evidence matrix}\label{tab:evidence-matrix}\\
\toprule
\textbf{Study} & \textbf{Primary task} & \textbf{Representative information used} & \textbf{Prioritization / output} & \textbf{Relevance to the present problem} \\
\midrule
\endfirsthead
\toprule
\textbf{Study} & \textbf{Primary task} & \textbf{Representative information used} & \textbf{Prioritization / output} & \textbf{Relevance to the present problem} \\
\midrule
\endhead
\bottomrule
\endlastfoot

\citet{iqbal2021} &
Culvert blockage classification &
Image-derived visual features &
Local blockage class &
Provides local condition information but not propagated network context \\
\addlinespace

\citet{vandaele2024} &
Trash-screen blockage monitoring &
Camera imagery and learned blockage features &
Asset-level blockage information for maintenance/flood management &
Demonstrates operational value of automated local disturbance information \\
\addlinespace

\citet{simone2022} &
Vulnerability and monitoring analysis &
Topology, intrinsic relevance, flow direction &
Tailored node/edge centrality ranking &
Shows that raw topology alone may not represent functional importance \\
\addlinespace

\citet{simone2023} &
Critical-pipe identification &
Connectivity and relevance-based centrality &
Pipe criticality ranking &
Supports network-aware prioritization but not current disturbance propagation \\
\addlinespace

\citet{dastgir2023} &
Critical-pipe analysis &
Runoff paths, edge betweenness, capacity-related information &
Critical-pipe ranking compared with hydrodynamic results &
Demonstrates lightweight graph screening and its accuracy/complexity tradeoff \\
\addlinespace

\citet{meijer2023} &
Weakest-link analysis &
Storage, discharge capacity, and graph flow paths &
Subsystem weakness / flood-performance screening &
Shows that prioritization may combine structure with operating-capacity information \\
\addlinespace

\citet{zhang2025} &
Siltation-chain risk analysis &
Pipeline risk factors, centrality, edge vulnerability, propagation paths &
Critical nodes and high-risk pathways &
Closest drainage precedent for propagated risk-path prioritization, but with siltation-specific semantics \\
\addlinespace

\citet{zhang2024} &
Hydraulic-state surrogate prediction &
Graph topology, recent hydraulic states, runoff, control/boundary information &
Predicted hydraulic variables rather than a direct risk ranking &
Defines the learned/high-fidelity boundary of the proposed lightweight risk-index model \\

\end{longtable}
}


The matrix shows that existing studies already provide local condition
signals, graph-based criticality, capacity-aware screening, propagation
pathways, and richer hydraulic state modeling. These strands remain
distributed across methods with different purposes and state semantics,
which creates an integrative question rather than an absence of ranking
methods \citep{iqbal2021,simone2022,dastgir2023,meijer2023,zhang2024,
zhang2025}.

\subsection{Parameter Patterns in Existing
Approaches}\label{parameter-patterns-in-existing-approaches}

An important model-design question concerns how many parameters should
be included in the generic core and its domain-specific instantiations.
The literature suggests that the
appropriate number cannot be fixed in advance; parameter choice follows
the semantics and purpose of the model \citep{simone2022,dastgir2023,
meijer2023,zhang2024,zhang2025}.

Drainage approaches already illustrate this diversity. Standard
centrality may use only graph structure, whereas tailored centrality can
incorporate intrinsic node relevance and flow direction \citep{simone2022},
\citep{simone2023}. \citet{dastgir2023} combine runoff-related path
information with capacity-related information for critical-pipe
analysis. \citet{meijer2023} focus on subsystem storage and
discharge capacities together with graph-based flow paths. Zhang et al.
\citep{zhang2025} combine pipeline risk factors, multiple centrality measures,
edge vulnerability, and propagation-path information. Hydraulic and
learned-surrogate approaches use substantially larger sets of rainfall,
boundary, geometric, control, and hydraulic-state variables \citep{zhang2024}.

These studies show that more parameters are not automatically better,
just as fewer parameters are not automatically preferable. Each added
variable should represent a distinct mechanism or provide empirically
useful domain information. Adding several correlated or weakly justified
parameters may reduce interpretability without improving the decision
task \citep{simone2022,dastgir2023,meijer2023,zhang2024}.

Across the reviewed methods, a defensible reusable core retains only
semantically distinct roles while allowing each domain to construct those
roles from several observed variables where justified. For example, a
drainage susceptibility mapping may combine hydraulic or physical
attributes without requiring a generic propagation equation to absorb
every raw parameter \citep{dastgir2023,meijer2023,zhang2025}.

Thus, the literature supports a generic core plus domain-supported
parameterization rather than a universal fixed list of physical
variables. The final number of drainage-specific parameters should be
determined from evidence, data availability, and validation rather than
from an arbitrary target count \citep{simone2022,dastgir2023,zhang2024}.


% \begin{quote}\textit{Table migration note: this comparison table was preserved separately in tables\_pending.tex and will be rebuilt cleanly in LaTeX before final submission.}\end{quote}

{\small
\renewcommand{\arraystretch}{1.15}
\begin{longtable}{@{}>{\raggedright\arraybackslash}p{0.22\dimexpr\linewidth-6\tabcolsep\relax}
                     >{\raggedright\arraybackslash}p{0.26\dimexpr\linewidth-6\tabcolsep\relax}
                     >{\raggedright\arraybackslash}p{0.22\dimexpr\linewidth-6\tabcolsep\relax}
                     >{\raggedright\arraybackslash}p{0.30\dimexpr\linewidth-6\tabcolsep\relax}@{}}
\caption{Parameter patterns in existing approaches}\label{tab:parameter-patterns}\\
\toprule
\textbf{Example approach} & \textbf{Representative information used} & \textbf{Role in model} & \textbf{Literature-derived implication} \\
\midrule
\endfirsthead
\toprule
\textbf{Example approach} & \textbf{Representative information used} & \textbf{Role in model} & \textbf{Literature-derived implication} \\
\midrule
\endhead
\bottomrule
\endlastfoot

Tailored centrality~\citep{simone2022,simone2023} &
Topology, intrinsic relevance, flow direction &
Improves structural ranking &
Topology can be enriched, but added variables need clear semantics \\
\addlinespace

Critical-pipe graph method~\citep{dastgir2023} &
Runoff/path information and capacity-related information &
Approximates criticality/flooding consequence &
Shows value of separating network paths from capacity response \\
\addlinespace

Weakest-link method~\citep{meijer2023} &
Storage, discharge capacity, flow paths &
Screens subsystem weakness &
Domain variables may be combined without redefining the generic graph concept \\
\addlinespace

Siltation-chain analysis~\citep{zhang2025} &
Pipeline risk, centrality, edge vulnerability, risk path &
Ranks critical nodes and paths &
Multiple parameters are justified when each represents a distinct mechanism \\
\addlinespace

GNN hydraulic surrogate~\citep{zhang2024} &
Hydraulic states, runoff, boundary, control variables, and graph topology &
Predicts physical hydraulic states &
High parameter count supports richer prediction but changes the task and data burden \\

\end{longtable}
}



\subsection{Reusable Architectures and Domain-Specific
Instantiation}\label{generic-core-and-domain-specific-instantiation}

A model can be generic in formulation without being empirically
generalizable to every domain. Network science repeatedly reuses
architectural ideas---nodes, directed relations, current state,
transmission, response, and aggregation---while redefining their meanings
for influence, epidemics, cascading infrastructure, cybersecurity,
finance, transportation, and drainage \citep{kempe2003,kermack1927,
buldyrev2010,sheyner2002,faraji2026,petrone2018,simone2022}.

This distinction is important because the word \emph{generalizable} can
imply demonstrated performance on unseen domains. A reusable architecture
supports cross-domain instantiation only when each application can
defensibly map its data and semantics to the abstract roles and validate
the resulting interpretation \citep{kavallieratos2021,faraji2026,
petrone2018}.

The core propagation logic can therefore remain fixed while a domain
supplies its own parameter definitions and, where justified, an
additional refinement layer. Drainage, cybersecurity, transportation,
and supply-chain applications require different parameter definitions,
evidence, and independent validation even when they share a graph-based
computational pattern \citep{kavallieratos2021,zhou2023,petrone2018,
dastgir2023}.

This structure resembles a reusable computational backbone with
application-specific adaptation, but it should not be interpreted as
machine-learning transfer learning. A learned surrogate transfers fitted
representations or parameters; an explicit architecture instead reuses a
declared computational rule whose variables must be reinterpreted and
validated in each domain \citep{zhang2024}.

\subsection{Synthesis of Research
Problems}\label{synthesis-of-research-problems}

The reviewed evidence supports three major research problems.

First, local condition does not by itself represent network-context
priority. Local monitoring can quantify condition at a component, but
comparative prioritization in an interconnected system requires
consideration of propagated network effects
\citep{iqbal2021,vandaele2024,simone2022,dastgir2023}.

Second, structural importance does not by itself represent current
operating response. Centrality and topology use the network as context,
but drainage and broader infrastructure studies repeatedly augment
structure with hydraulic relevance, capacity, storage/discharge
relationships, vulnerability, or other operating information
\citep{simone2022,dastgir2023,meijer2023,buldyrev2010,faraji2026}.

Third, existing propagation states are designed for different analytical
objectives. Activation, probability, infection, physical load, failure,
financial loss, and hydraulic state each imply different update and
aggregation rules. A bounded continuous non-probabilistic comparative
risk state therefore requires an explicitly declared semantic and
compatible propagation architecture \citep{kempe2003,kermack1927,
sheyner2002,buldyrev2010,faraji2026,zhang2024}.

Together, these problems create a need for a model that can integrate
local disturbance, current propagated state, and relation-specific
operating response while remaining interpretable enough to explain why
one node receives a higher comparative risk state than another
\citep{simone2022,dastgir2023,meijer2023,zhang2024}.

\subsection{Architectural and Integrative Research
Gap}\label{architectural-and-integrative-research-gap}

The literature does not support claiming that graph propagation,
recursive updating, capacity-based response, bounded state, edge
weighting, or multi-source aggregation are individually novel. Each has
precedent in established model families \citep{kempe2003,buldyrev2010,
munoz2017,faraji2026,petrone2018}.

The gap is instead architectural and integrative. Within the reviewed
drainage, infrastructure, cybersecurity, finance, transportation,
cascade, epidemic, influence-diffusion, and graph-surrogate evidence, no
clearly established lightweight directed-graph formulation was
identified that combines all of the following under one explicit
comparative-risk semantics: a bounded continuous non-probabilistic node
state; recursive propagation of current upstream risk; explicit
separation of directed reachability, optional transmission, and
relation-specific susceptibility; local disturbance at receiving nodes;
bounded aggregation of converging contributions; and explicit
recognition that shared-path or common-cause dependence is distinct from
local aggregation. This gap statement is a synthesis bounded to the
literature set examined in this chapter \citep{simone2022,dastgir2023,
meijer2023,zhang2025,kempe2003,buldyrev2010,sheyner2002}.

The reviewed evidence therefore establishes an architectural and
integrative research gap rather than the absence of individual graph,
propagation, or ranking mechanisms. A response to this gap would need to
combine a clearly defined bounded state, recursive updating,
relation-specific modifiers, receiving-node disturbance, convergence,
and explicit dependence limits while preserving the distinction between
generic computation and domain-specific evidence
\citep{kempe2003,buldyrev2010,sheyner2002,dastgir2023,zhang2025}.

Figure~\ref{fig:literature-gap} summarizes how the three reviewed
evidence streams lead to this architectural and integrative gap.

\begin{figure}[H]
\centering
\resizebox{0.85\textwidth}{!}{%
\begin{tikzpicture}[
  font=\sffamily,
  stream/.style={draw=gbgray, line width=.8pt, fill=white, minimum width=4.2cm, minimum height=1.45cm, align=center},
  gap/.style={draw=gbnavy, line width=1.1pt, fill=gblight, minimum width=7.8cm, minimum height=1.55cm, align=center},
  arr/.style={-{Latex[length=2.4mm]}, line width=.8pt, draw=gbgray}
]
\node[stream] (local) {\textbf{Local condition assessment}\\blockage, debris, inspection\\\footnotesize limitation: asset-level state};
\node[stream, right=8mm of local] (graph) {\textbf{Graph prioritization}\\centrality, paths, capacity\\\footnotesize limitation: structure $\neq$ current risk};
\node[stream, right=8mm of graph] (prop) {\textbf{Propagation families}\\activation, infection, failure\\\footnotesize limitation: different state semantics};
\node[gap, below=16mm of graph] (gap) {\textbf{Architectural and integrative gap}\\bounded comparative state + current-state recursion + separated mechanisms\\+ receiving-node disturbance + bounded convergence + explicit dependence limitation};
\draw[arr] (local.south) -- (gap.north -| local.south);
\draw[arr] (graph.south) -- (gap.north);
\draw[arr] (prop.south) -- (gap.north -| prop.south);
\end{tikzpicture}}
\caption{Literature-derived synthesis of existing approach families and the remaining architectural gap.}
\label{fig:literature-gap}
\end{figure}


\subsection{Chapter Synthesis}\label{chapter-synthesis}

The literature establishes a clear progression from practical drainage
prioritization to a broader network-risk modeling problem. Local
assessment provides necessary component-level condition information,
graph metrics provide structural context, and high-fidelity models
provide domain-rich physical behavior. Broader propagation research
contributes mechanisms for recursive updating, transmission,
susceptibility, aggregation, and dependence. However, these mechanisms
are distributed across models with different state semantics and
purposes \citep{kempe2003,kermack1927,buldyrev2010,sheyner2002,
kavallieratos2021,petrone2018}.

The drainage literature also demonstrates that prioritization methods
already range from local condition scoring to tailored centrality,
weakest-link analysis, siltation-chain risk, hydrodynamic consequence
analysis, and learned hydraulic surrogates. The research opportunity is
therefore not to introduce ranking itself, but to formulate a distinct
comparative-risk semantics and propagation process that can explain why
one node ranks above another after both local and network context are
considered \citep{iqbal2021,simone2022,dastgir2023,meijer2023,
zhang2024}.

Accordingly, the reviewed literature supports a methodological response
that declares its state semantics, separates topology from relation-level
response, propagates current state recursively, bounds converging
contributions, and limits its claims to comparative prioritization. The
specific formulation adopted by the study is presented in Chapter 3 so
that this chapter remains a synthesis of prior work rather than a
duplicate statement of the proposed model \citep{kempe2003,buldyrev2010,
munoz2017,dastgir2023,zhang2025}.
